\chapter{Tuberculosis Treatment}

\section{Overview}
TB treatment requires multiple antibiotics taken for 6-9 months. Successful treatment depends on taking all medications exactly as prescribed. The standard regimen includes four first-line drugs for the first two months, followed by two drugs for four additional months.

\section{Symptoms}
Treatment is indicated for:
\begin{itemize}
    \item Active pulmonary TB with positive sputum tests
    \item Active extrapulmonary TB confirmed by biopsy or culture
    \item Latent TB infection in high-risk individuals
    \item Drug-susceptible TB confirmed by sensitivity testing
\end{itemize}

\section{Causes}
Untreated TB can progress from latent infection to active disease, causing:
\begin{itemize}
    \item Lung destruction and cavitation
    \item Spread to kidneys, spine, brain, and bones
    \item Miliary TB (disseminated through bloodstream)
    \item TB meningitis (potentially fatal)
\end{itemize}

\section{Diagnosis}
Treatment selection depends on:
\begin{itemize}
    \item Drug susceptibility testing results
    \item HIV co-infection status
    \item Liver and kidney function
    \item Pregnancy status
    \item Drug interactions with other medications
\end{itemize}

\section{Treatment}
Standard first-line regimen:
\begin{itemize}
    \item \textbf{Intensive phase (2 months):} Isoniazid, Rifampin, Pyrazinamide, Ethambutol (HRZE)
    \item \textbf{Continuation phase (4 months):} Isoniazid and Rifampin (HR)
\end{itemize}

DOT (Directly Observed Therapy) ensures medication adherence. Side effects include hepatotoxicity, peripheral neuropathy, optic neuritis, and rash. Drug-resistant TB requires second-line agents like fluoroquinolones and injectable medications for 18-24 months.

\section{Prevention}
Preventing treatment failure and resistance requires:
\begin{itemize}
    \item Taking all medications exactly as prescribed
    \item Not skipping doses
    \item Completing the full treatment course
    \item Regular monitoring of liver function
    \item Reporting side effects immediately
\end{itemize}

\section{When to Seek Medical Care}
Seek care if experiencing:
\begin{itemize}
    \item Yellowing of skin or eyes (hepatotoxicity)
    \item Numbness or tingling in extremities
    \item Vision changes (orange discoloration of body fluids is expected with rifampin)
    \item Persistent nausea or vomiting
    \item Worsening symptoms despite treatment
\end{itemize}
